% Luc Destombe --- licence CC BY-NC-SA 4.0
% https://creativecommons.org/licenses/by-nc-sa/4.0/deed.fr
% Fichier autonome : compiler avec pdflatex, deux fois (signets, nombre de pages).
\def\lyceeversion{1.1}
\documentclass[a4paper,10pt]{article}
\makeatletter
\newif\iflycee@niveaux
\newif\iflycee@palatino
\lycee@palatinotrue

\RequirePackage[utf8]{inputenc}
\RequirePackage[T1]{fontenc}
\RequirePackage[french]{babel}
\RequirePackage{etoolbox}

\newcommand{\ShorthandsOff}{\@ifundefined{shorthandoff}{}{\shorthandoff{;:!?}}}
\newcommand{\ShorthandsOn}{\@ifundefined{shorthandon}{}{\shorthandon{;:!?}}}
\ShorthandsOff
\AtBeginDocument{\ShorthandsOn}
\AtBeginEnvironment{tikzpicture}{\ShorthandsOff}

\RequirePackage{geometry}
\RequirePackage{fancyhdr}
\RequirePackage{lastpage}
\RequirePackage{multicol}
\RequirePackage{array}
\RequirePackage{enumitem}
\RequirePackage{xcolor}
\RequirePackage{needspace}

\RequirePackage{amsmath,amssymb}
\RequirePackage{mathpazo}
\renewcommand{\ttdefault}{lmtt}
\DeclareMathSizes{10.95}{10.95}{8}{5.5}
\begingroup\catcode`\_=\active \catcode`\^=\active
\gdef_#1{\sb{\scriptscriptstyle #1}}
\gdef^#1{\sp{\scriptscriptstyle\lycee@moinsepais #1}}
\endgroup
\newcommand\lycee@moins{\mathbin{%
  \setbox\z@\hbox{$\m@th\scriptscriptstyle\lycee@moinsorig$}%
  \dimen@=.55\wd\z@ \advance\dimen@-.5pt
  \hbox to.75\wd\z@{\hss$\m@th\scriptscriptstyle\vcenter{\hbox{%
    \tikz\draw[line width=.5pt,line cap=round](0,0)--(\the\dimen@,0);}}$\hss}}}
\begingroup\lccode`\~=`\- \lowercase{\endgroup\let~}\lycee@moins
\newcommand\lycee@moinsepais{\mathcode`\-="8000 }
\AtBeginDocument{\mathchardef\lycee@moinsorig=\mathcode`\-\relax}
\AtBeginDocument{\catcode`\_=12 \mathcode`\_="8000
  \catcode`\^=12 \mathcode`\^="8000 }
\newcommand\lycee@exposants{%
  \fontdimen13\textfont2=.48\fontdimen6\textfont2
  \fontdimen14\textfont2=.48\fontdimen6\textfont2
  \fontdimen15\textfont2=.48\fontdimen6\textfont2
  \fontdimen13\scriptfont2=.48\fontdimen6\scriptfont2
  \fontdimen14\scriptfont2=.48\fontdimen6\scriptfont2
  \fontdimen15\scriptfont2=.48\fontdimen6\scriptfont2}
\every@math@size\expandafter{\the\every@math@size\lycee@exposants}
\newlength{\retraitcalculs}
\setlength{\retraitcalculs}{2em}

\RequirePackage{tikz}
\usetikzlibrary{arrows.meta,calc,babel,shapes.misc}
\RequirePackage[most]{tcolorbox}
\tcbuselibrary{skins,breakable}

\definecolor{coulDef}{HTML}{1F4E79}
\definecolor{coulProp}{HTML}{7030A0}
\definecolor{coulRem}{HTML}{C55A11}
\definecolor{coulEx}{HTML}{2E7D32}
\definecolor{coulExo}{HTML}{B03A2E}
\definecolor{coulPart}{HTML}{1F4E79}
\definecolor{coulMeth}{HTML}{1B6E4F}
\definecolor{coulCourbe}{HTML}{0B5ED7}
\definecolor{coulCourbeB}{HTML}{C00000}
\definecolor{coulCourbeC}{HTML}{2E7D32}
\definecolor{coulGrille}{HTML}{8C8C8C}
\definecolor{coulRep}{HTML}{1B6E4F}
\definecolor{coulBar}{HTML}{2E7D32}

\newcommand{\R}{\mathbb{R}}
\newcommand{\Rs}{\mathbb{R}^{*}}
\renewcommand{\le}{\leqslant}
\renewcommand{\ge}{\geqslant}
\renewcommand{\approx}{\simeq}
\newcommand{\vect}[1]{\overrightarrow{#1}}
\newcommand{\norme}[1]{\left\lVert #1 \right\rVert}
\newcommand{\intervalle}[2]{\left[#1\,;#2\right]}
\RequirePackage{cancel}
\renewcommand{\CancelColor}{\color{coulExo}}
\newcommand{\fois}[1]{\times{\color{coulExo}#1}}
\newcommand{\barre}[1]{\cancel{#1}}

\tikzset{grille/.style={coulGrille,line width=0.4pt}}
\newcommand{\quadrillage}[4]{\draw[grille,step=1] (#1,#2) grid (#3,#4);}
\tikzset{
  croix/.style={cross out,draw,line width=0.6pt,inner sep=0pt,minimum size=4pt},
  croixdroite/.style={croix,rotate=45,line width=0.8pt},
}
\newcommand{\pt}[3]{\path (#1) node[croix]{} node[#2,font=\small,inner sep=2.5pt]{$#3$};}

\newcommand{\lycee@repere}[4]{%
  \draw[grille,step=1] (#1,#2) grid (#3,#4);
  \draw[-{Stealth[length=2mm]},thick] (#1,0)--({#3+0.4},0) node[below left=-1pt]{$x$};
  \draw[-{Stealth[length=2mm]},thick] (0,#2)--(0,{#4+0.4}) node[below left=-1pt]{$y$};
  \node[below left,font=\scriptsize,inner sep=1.5pt] at (0,0) {$0$};}
\newcommand{\lycee@gradx}[1]{\foreach \k/\l in {#1}{%
  \draw (\k,0.07)--(\k,-0.07) node[below,font=\scriptsize,inner sep=1.5pt]{$\l$};}}
\newcommand{\lycee@grady}[1]{\foreach \k/\l in {#1}{%
  \draw (0.07,\k)--(-0.07,\k) node[left,font=\scriptsize,inner sep=1.5pt]{$\l$};}}
\newcommand{\lycee@intff}[2]{\left[#1\,;#2\right]}
\newcommand{\lycee@intfo}[2]{\left[#1\,;#2\right[}
\newcommand{\lycee@intof}[2]{\left]#1\,;#2\right]}
\newcommand{\lycee@intoo}[2]{\left]#1\,;#2\right[}
\newcommand{\lycee@ens}[1]{\left\{#1\right\}}
\AtBeginDocument{%
  \providecommand{\repere}{\lycee@repere}%
  \providecommand{\gradx}{\lycee@gradx}%
  \providecommand{\grady}{\lycee@grady}%
  \providecommand{\Cf}{\mathcal{C}\sb{\scriptscriptstyle f}}%
  \providecommand{\Cg}{\mathcal{C}\sb{\scriptscriptstyle g}}%
  \providecommand{\intff}{\lycee@intff}%
  \providecommand{\intfo}{\lycee@intfo}%
  \providecommand{\intof}{\lycee@intof}%
  \providecommand{\intoo}{\lycee@intoo}%
  \providecommand{\ens}{\lycee@ens}}

\newlist{questions}{enumerate}{3}
\setlist[questions,1]{label=\arabic*\textdegree),leftmargin=*,itemsep=2pt,topsep=3pt}
\setlist[questions,2]{label=\alph*),leftmargin=*,itemsep=1pt,topsep=2pt}
\setlist[questions,3]{label=\roman*),leftmargin=*,itemsep=1pt,topsep=1pt}
\setlist[itemize]{label=\textbullet,leftmargin=*,itemsep=2pt,topsep=3pt}

\newcommand{\besoinplace}[1]{\par\penalty\@M\begingroup
  \dimen@=#1\relax\dimen@ii=\pagegoal\advance\dimen@ii-\pagetotal
  \ifdim\dimen@>\dimen@ii\ifdim\dimen@ii>\z@\vfil\fi\break\fi\endgroup}

\newcommand{\lycee@pied}{\thepage/\pageref{LastPage}}
\newcommand{\entete}[2]{%
  \pagestyle{fancy}\fancyhf{}%
  \fancyhead[L]{\footnotesize\color{coulPart}\textbf{#1}}%
  \fancyhead[R]{\footnotesize\color{coulPart}\textbf{#2}}%
  \fancyfoot[C]{\footnotesize\lycee@pied}%
  \renewcommand{\headrulewidth}{0.4pt}%
  \renewcommand{\headrule}{\hbox to\headwidth{%
    \color{coulPart!40}\leaders\hrule height \headrulewidth\hfill}}}

\RequirePackage{ccicons}
\newcommand{\lycee@licence}{{\scriptsize\color{gray}%
  \copyright\ Luc Destombe --- \ccbyncsa\ CC BY-NC-SA 4.0}}
\newif\iflycee@licence \lycee@licencetrue
\newcommand{\sanslicence}{\lycee@licencefalse}
\AtBeginDocument{\iflycee@licence\AddToHookNext{shipout/foreground}{%
  \put(0,0){\rlap{\kern\dimexpr1in+\hoffset+\oddsidemargin+\textwidth\relax
    \llap{\raisebox{-\dimexpr1in+\voffset+\topmargin+\headheight+\headsep
      +\textheight+\footskip\relax}{\lycee@licence}}}}}\fi}

\newcommand{\lycee@site}{faire-des-maths.fr}
\newcommand{\lycee@version}{{\scriptsize\color{gray}\lycee@site
  \ifdefined\lyceeversion\ --- Version~\lyceeversion\fi}}
\AtBeginDocument{\iflycee@licence\AddToHookNext{shipout/foreground}{%
  \put(0,0){\rlap{\kern\dimexpr1in+\hoffset+\oddsidemargin\relax
    \raisebox{-\dimexpr1in+\voffset+\topmargin+\headheight+\headsep
      +\textheight+\footskip\relax}{\lycee@version}}}}\fi}

\newcommand{\formulesserrees}{%
  \setlength{\abovedisplayskip}{3pt}\setlength{\belowdisplayskip}{3pt}%
  \setlength{\abovedisplayshortskip}{2pt}\setlength{\belowdisplayshortskip}{3pt}}

\setlength{\parindent}{0pt}

\RequirePackage[implicit=false,hidelinks,pdfencoding=auto,psdextra,bookmarksopen,
  bookmarksopenlevel=1,pdfcreator={},pdfproducer={}]{hyperref}
\RequirePackage{bookmark}
\hypersetup{pdfauthor={Luc Destombe},pdfkeywords={CC BY-NC-SA 4.0}}
\newcounter{lycee@signet}
\newcommand{\signet}[3][]{\stepcounter{lycee@signet}%
  \edef\lycee@dest{\if\relax\detokenize{#1}\relax
    signet.\arabic{lycee@signet}\else#1\fi}%
  \hypertarget{\lycee@dest}{}\bookmark[level=#2,dest=\lycee@dest]{#3}}
\pdfstringdefDisableCommands{%
  \def\R{R}\def\vect#1{#1}\def\Cf{Cf}\def\Cg{Cg}\def\fois#1{×#1}%
  \def\barre#1{#1}\def\intervalle#1#2{[#1;#2]}\def\intff#1#2{[#1;#2]}%
  \def\textsuperscript#1{#1}\def\dfrac#1#2{#1/#2}\def\frac#1#2{#1/#2}%
  \def\vec#1{#1⃗}}%
\RequirePackage[official]{eurosym}

\geometry{top=0.8cm,bottom=1.3cm,left=1cm,right=1cm,footskip=0.6cm}

\pagestyle{fancy}
\fancyhf{}
\renewcommand{\headrulewidth}{0pt}
\fancyfoot[C]{\footnotesize Page \thepage{} / \pageref{LastPage}}

\newcommand{\titrefiche}[2]{%
  \begin{center}
    \Large\textbf{#1}\\[2pt]
    \large\textbf{#2}\\[2pt]
  \end{center}
  \vspace{0.10cm}}

\setlength{\columnseprule}{0.4pt}
\setlength{\columnsep}{15pt}
\renewcommand{\columnseprulecolor}{\color{black!45}}
\setlength{\multicolsep}{2pt}

\newcounter{partiefiche}
\renewcommand{\thepartiefiche}{\Roman{partiefiche}}
\newif\ifapresbandeau
\newcommand{\lycee@bandeau}[1]{%
  \par\penalty-600
  \vspace{0.08cm}%
  \noindent\colorbox{coulPart}{%
    \parbox{\dimexpr\linewidth-2\fboxsep\relax}{%
      \color{white}\bfseries\raggedright\strut#1\strut}}%
  \nopagebreak\par\nopagebreak
  \vspace{0.06cm}\nopagebreak
  \apresbandeautrue}
\newcommand{\partiefiche}[1]{%
  \refstepcounter{partiefiche}%
  \lycee@bandeau{\signet{1}{\thepartiefiche{} --- #1}\thepartiefiche{} --- #1}}
\newcommand{\partiefichesansnum}[1]{\lycee@bandeau{\signet{1}{#1}#1}}

\newcounter{exo}
\newlength{\espaceexo}\setlength{\espaceexo}{7pt}
\newlength{\espacetitre}\setlength{\espacetitre}{2pt}
\newcommand{\exo}[1][\relax]{%
  \par
  \ifapresbandeau\apresbandeaufalse\else\penalty-150\addvspace{\espaceexo}\fi
  \refstepcounter{exo}%
  \noindent\signet[exercice-\arabic{exo}]{2}{Exercice \arabic{exo}}%
  \textbf{\color{coulExo}\underline{Exercice \theexo}}%
  \ifx\relax#1\relax\else\textbf{ : #1}\fi
  \par\nopagebreak\vspace{\espacetitre}\nopagebreak
  \ignorespaces}

\setlist[enumerate]{nosep,leftmargin=*,itemsep=2pt,topsep=2pt}
\setlist[itemize]{nosep,leftmargin=*,topsep=2pt}
\newcommand{\choix}[3]{\par\hspace*{1em}%
  \makebox[0.3\linewidth][l]{a) #1}\makebox[0.3\linewidth][l]{b) #2}c) #3}
\newcommand{\deux}[2]{\par\hspace*{0.6em}%
  \makebox[0.48\linewidth][l]{#1}#2}
\newcommand{\trois}[3]{\par\hspace*{0.6em}%
  \makebox[0.32\linewidth][l]{#1}\makebox[0.32\linewidth][l]{#2}#3}

\renewenvironment{center}{\par\addvspace{3pt}\centering}{\par\addvspace{3pt}}
\formulesserrees
\AtBeginDocument{\formulesserrees}
\clubpenalty=10000
\widowpenalty=10000
\displaywidowpenalty=10000
\brokenpenalty=10000
\setlength{\parskip}{0pt}

\RequirePackage{ragged2e}
\setlength{\RaggedRightRightskip}{0pt plus 1fil}
\AtBeginDocument{\RaggedRight\binoppenalty=10000 \relpenalty=10000 }
\g@addto@macro\@parboxrestore{\RaggedRight}
\makeatother
\geometry{bottom=1cm}

\setlength{\espaceexo}{4pt}
\setlist[enumerate]{nosep,leftmargin=*,itemsep=1pt,topsep=1pt}

\ShorthandsOff
\tikzset{
  droite/.style={coulCourbe,line width=1pt},
  droiteB/.style={coulCourbeB,line width=1pt},
  droiteC/.style={coulCourbeC,line width=1pt},
}
\ShorthandsOn
\newcommand{\coord}[2]{\begin{pmatrix}#1\\#2\end{pmatrix}}

\begin{document}

\titrefiche{Seconde --- Équations de droites}{Fiche d'exercices du chapitre}

{\small\itshape Le repère est orthonormé dans toute la fiche ; on trace les repères sur
le cahier, un carreau pour unité.\par}
\vspace{0.12cm}

\begin{multicols}{2}\raggedcolumns

\partiefiche{Vecteur directeur d'une droite}

\exo[Lire un vecteur directeur]
\begin{center}
\begin{tikzpicture}[x=0.42cm,y=0.42cm]
  \repere{-4}{-3}{5}{4}
  \gradx{-3,-2,-1,1,2,3,4} \grady{-2,-1,1,2,3}
  \draw[droite] (-4,-2.75) -- (5,4);
  \draw[droiteB] (-4,4) -- (5,-0.5);
  \draw[droiteC] (0,-3) -- (1.75,4);
  \node[coulCourbe,font=\footnotesize,left] at (4.4,4) {$(d_{1})$};
  \node[coulCourbeB,font=\footnotesize,above] at (-3,3.6) {$(d_{2})$};
  \node[coulCourbeC,font=\footnotesize,right] at (1.7,3.4) {$(d_{3})$};
\end{tikzpicture}
\end{center}
Pour chaque droite $(d_{1})$, $(d_{2})$ et $(d_{3})$, marquer d'une croix deux nœuds du
quadrillage par lesquels elle passe, puis lire un vecteur directeur.

\exo[Tracer une droite]
Dans un repère, tracer chaque droite, puis donner les coordonnées d'un second point de
cette droite.
\begin{enumerate}[label=\alph*)]
  \item $(d_{1})$ passe par $A(-3\,;1)$, de vecteur directeur $\vec{u}\coord{3}{1}$.
  \item $(d_{2})$ passe par $B(-1\,;4)$, de vecteur directeur $\vec{v}\coord{1}{-3}$.
  \item $(d_{3})$ passe par $C(-2\,;-2)$, de vecteur directeur $\vec{w}\coord{2}{0}$.
\end{enumerate}

\exo[Est-ce un vecteur directeur ?]
\begin{enumerate}[label=\alph*)]
  \item On considère les points $A(2\,;5)$ et $B(6\,;-1)$. Les vecteurs
        $\vec{u}\coord{-2}{3}$ et $\vec{v}\coord{6}{4}$ sont-ils des vecteurs directeurs
        de la droite $(AB)$ ?
  \item On considère les points $C(3\,;-2)$ et $D(1\,;4)$. Calculer le nombre réel $x$
        tel que $\vec{w}\coord{1}{x}$ soit un vecteur directeur de $(CD)$.
\end{enumerate}

\exo[Un point sur la droite ?]
La droite $(d)$ passe par $A(1\,;-1)$ et a pour vecteur directeur $\vec{u}\coord{2}{3}$.
Pour chaque point $M(5\,;5)$ et $N(-3\,;-5)$, dire s'il est sur $(d)$, en calculant un
déterminant.

\partiefiche{Équation cartésienne d'une droite}

\exo[Exploiter une équation]
La droite $(d)$ a pour équation $2x-3y+6=0$.
\begin{enumerate}[label=\alph*)]
  \item Donner un vecteur directeur de $(d)$.
  \item Pour chaque point $A(0\,;2)$, $B(3\,;4)$ et $C(1\,;3)$, dire s'il est sur $(d)$.
  \item Calculer l'ordonnée du point de $(d)$ d'abscisse $1$.
  \item Calculer l'abscisse du point de $(d)$ d'ordonnée $0$.
  \item Tracer la droite $(d)$.
\end{enumerate}

\exo[Des vecteurs directeurs]
Pour chaque droite, donner un vecteur directeur (on écrit d'abord l'équation sous la
forme $ax+by+c=0$ si besoin).
\deux{a)~$3x-5y+2=0$}{b)~$4x+3y-5=0$}
\deux{c)~$-x+6y+2=0$}{d)~$2x=3y$}

\exo[Appartenir à une droite]
La droite $(d)$ a pour équation $x+3y-6=0$.
\begin{enumerate}[label=\alph*)]
  \item Pour chaque point $E(0\,;2)$, $F(3\,;1)$, $G(-3\,;3)$ et $H(2\,;2)$, écrire une
        phrase : « Le point … appartient (ou n'appartient pas) à la droite $(d)$ car … ».
  \item Le point $K(k\,;0)$ est sur $(d)$. Calculer $k$.
  \item Le point $L(1\,;y)$ est sur $(d)$. Calculer $y$.
\end{enumerate}

\exo[Avec le déterminant]
La droite $(d)$ passe par $A(2\,;1)$ et a pour vecteur directeur $\vec{u}\coord{1}{3}$.
On considère un point $M(x\,;y)$.
\begin{enumerate}[label=\alph*)]
  \item Calculer les coordonnées de $\vect{AM}$ en fonction de $x$ et $y$.
  \item Calculer $\det(\vect{AM}\,;\vec{u})$ et l'écrire sous forme réduite.
  \item En déduire une équation cartésienne de $(d)$, et vérifier que $\vec{u}$ est bien
        de la forme $\coord{-b}{a}$.
\end{enumerate}

\partiefiche{Trouver une équation cartésienne}

\exo[Un point et un vecteur directeur]
Pour chaque droite, trouver une équation cartésienne.
\begin{enumerate}[label=\alph*)]
  \item $(d_{1})$ passe par $A(2\,;-3)$, de vecteur directeur $\vec{u}\coord{-1}{4}$.
  \item $(d_{2})$ passe par $B(-1\,;2)$, de vecteur directeur $\vec{v}\coord{3}{2}$.
  \item $(d_{3})$ passe par $C(0\,;-4)$, de vecteur directeur $\vec{w}\coord{5}{-2}$.
\end{enumerate}

\exo[Deux points]
Pour chaque droite, trouver une équation cartésienne, puis vérifier avec le second point.
\deux{a)~$(AB)$ : $A(1\,;3)$ et $B(3\,;7)$}{b)~$(CD)$ : $C(-2\,;3)$ et $D(4\,;0)$}

\exo[Les droites d'un triangle]
On considère les points $A(-1\,;1)$, $B(3\,;3)$ et $C(1\,;-3)$.
\begin{enumerate}[label=\alph*)]
  \item Trouver une équation cartésienne de chaque droite $(AB)$ et $(BC)$.
  \item Calculer les coordonnées du milieu $I$ de $[BC]$.
  \item Trouver une équation cartésienne de la médiane $(AI)$.
\end{enumerate}

\exo[Du graphique à l'équation]
\begin{center}
\begin{tikzpicture}[x=0.42cm,y=0.42cm]
  \repere{-4}{-3}{5}{4}
  \gradx{-3,-2,-1,1,2,3,4} \grady{-2,-1,1,2,3}
  \draw[droite] (-4,-0.5) -- (5,4);
  \draw[droiteB] (-0.5,4) -- (3,-3);
  \node[coulCourbe,font=\footnotesize,left] at (4.3,4) {$(d_{1})$};
  \node[coulCourbeB,font=\footnotesize,right] at (2.6,-2.6) {$(d_{2})$};
\end{tikzpicture}
\end{center}
Pour chaque droite $(d_{1})$ et $(d_{2})$, marquer deux nœuds, lire un vecteur directeur,
puis trouver une équation cartésienne.

\partiefiche{Équation réduite ; passer de l'une à l'autre}

\exo[De la cartésienne à la réduite]
Pour chaque droite, trouver l'équation réduite, puis donner sa pente et son ordonnée à
l'origine.
\deux{a)~$2x+y-5=0$}{b)~$6x-3y+9=0$}
\deux{c)~$x+3y-6=0$}{d)~$4x-6y+3=0$}

\exo[De la réduite à la cartésienne]
Pour chaque droite, trouver une équation cartésienne à coefficients entiers.
\deux{a)~$y=3x-2$}{b)~$y=-x+4$}
\deux{c)~$y=\frac{2}{3}x-1$}{d)~$y=-\frac{1}{4}x+2$}

\exo[Pente et vecteur directeur]
Pour chaque droite $(d_{1})$ : $y=2x+1$, $(d_{2})$ : $y=-3x$ et
$(d_{3})$ : $y=\frac{1}{3}x+2$, donner un vecteur directeur (à coordonnées entières),
puis tracer la droite à partir de son ordonnée à l'origine.

\exo[Par un point et la pente, par deux points]
Pour chaque droite, trouver l'équation réduite (chapitre 5, méthode 8).
\begin{enumerate}[label=\alph*)]
  \item La droite passe par $A(-1\,;4)$ et a pour pente $2$.
  \item La droite $(BC)$, avec $B(1\,;-1)$ et $C(3\,;5)$.
  \item La droite $(DE)$, avec $D(-2\,;3)$ et $E(4\,;0)$.
\end{enumerate}

\exo[Le prix d'une course]
Dans un repère où $x$ est la distance parcourue (en km) et $y$ le prix d'une course de
taxi (en euros), les points qui représentent les courses sont sur la droite d'équation
$4x-5y+20=0$.
\begin{enumerate}[label=\alph*)]
  \item Écrire l'équation sous sa forme réduite.
  \item Que représentent la pente et l'ordonnée à l'origine pour le client ?
  \item Calculer le prix d'une course de $15$~km.
  \item Quelle distance peut-on parcourir pour $28$~€ ?
\end{enumerate}

\partiefiche{Droites parallèles}

\exo[Parallèles ou sécantes]
Pour chaque couple de droites, dire si les droites sont parallèles ou sécantes.
\begin{enumerate}[label=\alph*)]
  \item $(d_{1})$ : $3x-2y+1=0$ et $(d_{2})$ : $-6x+4y+7=0$.
  \item $(d_{3})$ : $x+y-2=0$ et $(d_{4})$ : $x-y+3=0$.
  \item $(d_{5})$ : $2x+5y=0$ et $(d_{6})$ : $4x+10y-1=0$.
  \item $(d_{7})$ : $x-4y+2=0$ et $(d_{8})$ : $\frac{1}{4}x-y+5=0$.
\end{enumerate}

\exo[La forme la plus pratique]
Pour chaque couple de droites, écrire les deux équations sous la même forme (la plus
pratique), puis dire si les droites sont parallèles ou sécantes.
\begin{enumerate}[label=\alph*)]
  \item $(d_{1})$ : $y=2x-1$ et $(d_{2})$ : $4x-2y+3=0$.
  \item $(d_{3})$ : $y=-x+2$ et $(d_{4})$ : $x-y=0$.
\end{enumerate}

\exo[La parallèle qui passe par un point]
Pour chaque cas, trouver une équation cartésienne de la parallèle.
\begin{enumerate}[label=\alph*)]
  \item La parallèle à $(d)$ : $2x-y+4=0$ qui passe par $A(3\,;1)$.
  \item La parallèle à $(d')$ : $x+3y-1=0$ qui passe par $B(-2\,;4)$.
  \item La parallèle à $(EF)$ qui passe par $G(1\,;-2)$, avec $E(0\,;1)$ et $F(2\,;4)$.
\end{enumerate}

\exo[Un trapèze]
On considère les points $A(-2\,;-1)$, $B(4\,;1)$, $C(3\,;4)$ et $D(0\,;3)$.
\begin{enumerate}[label=\alph*)]
  \item Démontrer que les droites $(AB)$ et $(DC)$ sont parallèles.
  \item Les droites $(AD)$ et $(BC)$ sont-elles parallèles ? Quelle est la nature de
        $ABCD$ ?
  \item Trouver une équation cartésienne de $(AB)$, puis de la parallèle à $(AB)$ qui
        passe par $D$. Vérifier que $C$ est sur cette parallèle.
\end{enumerate}

\partiefiche{Droites horizontales et verticales}

\exo[Tracer]
Pour chaque droite, dire si elle est horizontale ou verticale, puis la tracer dans un
même repère.
\deux{a)~$(d_{1})$ : $3y-6=0$}{b)~$(d_{2})$ : $x+3=0$}
\deux{c)~$(d_{3})$ : $2y+4=0$}{d)~$(d_{4})$ : $4x-2=0$}

\exo[Trouver une équation]
Pour chaque droite, trouver une équation cartésienne, puis la simplifier en $x=k$ ou
$y=k$.
\begin{enumerate}[label=\alph*)]
  \item La droite $(AB)$, avec $A(2\,;-1)$ et $B(2\,;5)$.
  \item La droite $(CD)$, avec $C(-3\,;4)$ et $D(5\,;4)$.
  \item La droite qui passe par $E(-1\,;3)$, de vecteur directeur $\vec{\imath}\coord{1}{0}$.
\end{enumerate}

\exo[Les côtés d'un rectangle]
On considère le rectangle $ABCD$, avec $A(-2\,;-1)$, $B(3\,;-1)$, $C(3\,;2)$ et
$D(-2\,;2)$.
\begin{enumerate}[label=\alph*)]
  \item Donner une équation de chaque droite $(AB)$, $(BC)$, $(CD)$ et $(DA)$.
  \item Pour la droite $(BC)$, écrire une phrase qui dit à quel axe elle est parallèle.
  \item Donner une équation de chacun des deux axes de symétrie du rectangle.
\end{enumerate}

\partiefiche{Point d'intersection de deux droites sécantes}

\exo[Calculer un point d'intersection]
Pour chaque couple de droites, calculer les coordonnées du point d'intersection.
\begin{enumerate}[label=\alph*)]
  \item $(d_{1})$ : $y=3x-1$ et $(d_{2})$ : $y=-x+7$.
  \item $(d_{3})$ : $y=2x+3$ et $(d_{4})$ : $y=-x+1$.
  \item $(d_{5})$ : $y=\frac{1}{2}x+1$ et $(d_{6})$ : $y=-x+4$.
\end{enumerate}

\exo[D'abord réduire]
Pour chaque couple de droites, écrire les équations sous la forme la plus pratique, puis
calculer les coordonnées du point d'intersection.
\begin{enumerate}[label=\alph*)]
  \item $(d_{1})$ : $2x-y+1=0$ et $(d_{2})$ : $x+y-4=0$.
  \item $(d_{3})$ : $x+2=0$ et $(d_{4})$ : $y=3x+5$.
\end{enumerate}

\exo[Louer un vélo]
Une boutique loue des vélos à deux tarifs : tarif A, $6$~€ de frais d'inscription
puis $1{,}50$~€ par heure ; tarif B, $3$~€ par heure. On note $x$ la durée de location
(en heures) et $y$ le prix (en euros).
\begin{enumerate}[label=\alph*)]
  \item Écrire l'équation réduite de la droite de chaque tarif.
  \item Calculer les coordonnées du point d'intersection des deux droites.
  \item Quel tarif conseiller pour une location de $3$~h ? de $6$~h ?
\end{enumerate}

\exo[Deux droites sécantes]
On considère les points $A(0\,;1)$, $B(2\,;5)$, $C(-1\,;6)$ et $D(5\,;0)$.
\begin{enumerate}[label=\alph*)]
  \item Trouver les équations réduites des droites $(AB)$ et $(CD)$.
  \item Justifier que ces droites sont sécantes.
  \item Calculer les coordonnées de leur point d'intersection.
\end{enumerate}

\partiefiche{Problèmes de géométrie repérée}

\exo[Dans un carré]
$ABCD$ est un carré de côté $4$~cm ; $I$ est le milieu de $[BC]$ et $J$ celui de $[CD]$.
On se place dans le repère orthonormé d'origine $A$, d'axes $(AB)$ et $(AD)$, d'unité
$1$~cm.
\begin{enumerate}[label=\alph*)]
  \item Donner les coordonnées des points $A$, $B$, $C$, $D$, $I$ et $J$.
  \item Trouver les équations réduites des droites $(AI)$ et $(BJ)$.
  \item Calculer les coordonnées du point $K$, intersection de $(AI)$ et $(BJ)$.
  \item Démontrer que $DK=DA$.
\end{enumerate}

\exo[Trois points alignés]
$ABCD$ est un rectangle avec $AB=6$~cm et $AD=4$~cm ; $K$ est le milieu de $[AB]$ et $M$
le point de $[AC]$ tel que $\vect{AM}=\frac{1}{3}\vect{AC}$. On se place dans le repère
orthonormé d'origine $A$, d'axes $(AB)$ et $(AD)$, d'unité $1$~cm.
\begin{enumerate}[label=\alph*)]
  \item Donner les coordonnées des points $A$, $B$, $C$, $D$ et $K$.
  \item Calculer les coordonnées du point $M$.
  \item Trouver une équation cartésienne de la droite $(DK)$.
  \item En déduire que les points $D$, $M$ et $K$ sont alignés.
\end{enumerate}

\exo[Les médianes d'un triangle]
On considère les points $A(0\,;0)$, $B(6\,;0)$ et $C(0\,;6)$ ; $I$, $J$ et $K$ sont les
milieux de $[BC]$, $[AC]$ et $[AB]$.
\begin{enumerate}[label=\alph*)]
  \item Calculer les coordonnées de $I$, $J$ et $K$.
  \item Trouver les équations réduites des médianes $(AI)$ et $(BJ)$.
  \item Calculer les coordonnées de leur point d'intersection $G$.
  \item Démontrer que $G$ est aussi sur la troisième médiane $(CK)$.
\end{enumerate}

\end{multicols}

\end{document}
